\section{评测体系升级：Towbench 与工业级仿真}

讲者展示了 Towbench（Tool-Agent-User benchmark）的设计目标：把 Agent 放入包含工具、策略、多轮交互、动态状态变化的近真实环境中评测，而不是只做静态问答。

\subsection{为什么 unit test 不够}

Agent 是非确定性系统。相同输入在不同上下文、不同系统状态下可能产生不同动作。传统 unit test 能覆盖函数正确性，但难覆盖多轮对话、工具副作用、策略约束与用户情绪扰动。

\begin{knowledgebox}{Towbench 的四类必要构件}
\begin{itemize}[leftmargin=1.5em]
    \item 真实域任务：如电信、零售、航旅客服；
    \item 可执行工具：数据库与 API，支持状态变更；
    \item 用户模拟器：含 persona、情绪和噪声；
    \item 成功判定：基于可验证动作，而非仅 LLM-as-judge。
\end{itemize}
\end{knowledgebox}

\subsection{pass@k 与规模稳定性}

讲者强调“做对一次”不等于“持续做对”。在百万级会话里，哪怕单次成功率很高，重复调用后的失败暴露也会急剧放大。因而需要关注跨多轮、多样本的稳定指标，如 \texttt{pass@k} 的业务化版本。

\begin{importantbox}{评测目标的转移}
研究阶段常问“最优样例能到多高”；生产阶段必须问“在复杂分布上能稳定到什么水平”。Towbench 和工业仿真把这件事前置到上线前，而不是等线上事故驱动。
\end{importantbox}

\subsection{本章小结}

评测不是发布前打勾项，而是持续运营机制。Agent 团队需要把 benchmark、simulation、回归测试和线上回放连接成闭环，才能持续提升真实解决率。

